\begin{statementbox}
\textbf{条件：}设 $p\ne0$，$O=(0,0)$ 是抛物线的顶点，$A,B$ 为曲线上\textbf{两个不同且均不为 $O$} 的点。讨论三角形 $OAB$ 的三条边所在直线的斜率。

\begin{center}
% p=1. Right panel is the exact x/y interchange of left panel.
\begin{tikzpicture}[font=\small,>=stealth,line join=round]
\draw[gray!45,rounded corners=3pt] (0,0) rectangle (14.2,10.5);
\node at (7.1,9.95) {看同一个三角形的三条边：蓝 $OA$，橙 $OB$，绿 $AB$};
\node at (3.55,9.05) {$y^2=2x$：横向开口};
\node at (10.65,9.05) {$x^2=2y$：纵向开口};
\begin{scope}[shift={(2.0,5.0)},scale=1.05]
\draw[->,gray] (-1.35,0)--(3.0,0) node[above] {$x$};
\draw[->,gray] (0,-1.65)--(0,3.0) node[left] {$y$};
\draw[teal!65!black,thick,samples=100,domain=-1.45:2.3,variable=\t] plot ({\t*\t/2},\t);
\fill[blue!4] (0,0)--(2,2)--(.5,-1)--cycle;
\draw[blue!75!black,very thick] (0,0)--(2,2);
\draw[orange!85!black,very thick] (0,0)--(.5,-1);
\draw[green!45!black,very thick] (2,2)--(.5,-1);
\fill (0,0) circle(2pt) node[below left=7pt] {$O$};
\fill (2,2) circle(2pt) node[below right=8pt] {$A$};
\fill (.5,-1) circle(2pt) node[right=8pt] {$B$};
\end{scope}
\begin{scope}[shift={(9.1,5.0)},scale=1.05]
\draw[->,gray] (-1.65,0)--(3.0,0) node[above] {$x$};
\draw[->,gray] (0,-1.35)--(0,3.0) node[left] {$y$};
\draw[teal!65!black,thick,samples=100,domain=-1.45:2.3,variable=\t] plot (\t,{\t*\t/2});
\fill[blue!4] (0,0)--(2,2)--(-1,.5)--cycle;
\draw[blue!75!black,very thick] (0,0)--(2,2);
\draw[orange!85!black,very thick] (0,0)--(-1,.5);
\draw[green!45!black,very thick] (2,2)--(-1,.5);
\fill (0,0) circle(2pt) node[below left=7pt] {$O$};
\fill (2,2) circle(2pt) node[below right=8pt] {$A$};
\fill (-1,.5) circle(2pt) node[above=14pt] {$B$};
\end{scope}
\node at (3.55,2.65) {$\textcolor{blue!75!black}{k_{OA}=1},\quad\textcolor{orange!85!black}{k_{OB}=-2},\quad\textcolor{green!45!black}{k_{AB}=2}$};
\node at (10.65,2.65) {$\textcolor{blue!75!black}{k_{OA}=1},\quad\textcolor{orange!85!black}{k_{OB}=-\frac12},\quad\textcolor{green!45!black}{k_{AB}=\frac12}$};
\node at (3.55,1.55) {$\displaystyle\frac11+\frac1{-2}=\frac12$};
\node at (10.65,1.55) {$\displaystyle 1+\left(-\frac12\right)=\frac12$};
\node at (7.1,.55) {两图同取 $p=1$，坐标互换且单位长度相同；三条边的颜色一一对应。};
\end{tikzpicture}

\end{center}

\textbf{结论一：横向开口}（$p>0$ 向右，$p<0$ 向左）。对于
\[
C_1:\ y^2=2px,
\]
当 $k_{OA},k_{OB},k_{AB}$ 均存在时，
\[
\boxed{\frac1{k_{OA}}+\frac1{k_{OB}}=\frac1{k_{AB}}}.
\tag{H}
\]

\textbf{结论二：纵向开口}（$p>0$ 向上，$p<0$ 向下）。对于
\[
C_2:\ x^2=2py,
\]
总有
\[
\boxed{k_{OA}+k_{OB}=k_{AB}}.
\tag{V}
\]
这里的 $k_{AB}=0$ \textbf{允许}，不需要每个斜率都非零。

\textbf{一句话记忆：}横向开口\textbf{倒斜率相加}；纵向开口\textbf{斜率相加}。
\end{statementbox}

\clearpage
\begin{statementbox}
\textbf{不怕竖直弦的加强式。}设 $A=(x_A,y_A),B=(x_B,y_B)$，对于 $y^2=2px$，不同点的纵坐标必不同。只要 $A,B\ne O$，恒有
\begin{center}% p=1, free coordinates +1,-1, reflected with identical scale.
\begin{tikzpicture}[font=\small,>=stealth]
\draw[gray!45,rounded corners=3pt] (0,0) rectangle (14.2,10.0);
\node at (7.1,9.45) {对称的两个端点：竖直弦与水平弦的斜率含义不同};
\node at (3.55,8.45) {$y^2=2x$};
\node at (10.65,8.45) {$x^2=2y$};
\begin{scope}[shift={(2.6,5.1)},scale=1.25]
\draw[->,gray] (-1.5,0)--(2.1,0) node[above] {$x$};
\draw[->,gray] (0,-1.5)--(0,2.1) node[left] {$y$};
\draw[teal!65!black,thick,samples=90,domain=-1.4:1.8,variable=\t] plot ({\t*\t/2},\t);
\draw[blue!75!black,thick] (0,0)--(.5,1);
\draw[orange!85!black,thick] (0,0)--(.5,-1);
\draw[green!45!black,very thick] (.5,1)--(.5,-1);
\fill (0,0) circle(2pt) node[below left=7pt] {$O$};
\fill (.5,1) circle(2pt) node[right=9pt] {$A$};
\fill (.5,-1) circle(2pt) node[right=9pt] {$B$};
\end{scope}
\begin{scope}[shift={(9.7,5.1)},scale=1.25]
\draw[->,gray] (-1.5,0)--(2.1,0) node[above] {$x$};
\draw[->,gray] (0,-1.5)--(0,2.1) node[left] {$y$};
\draw[teal!65!black,thick,samples=90,domain=-1.4:1.8,variable=\t] plot (\t,{\t*\t/2});
\draw[blue!75!black,thick] (0,0)--(1,.5);
\draw[orange!85!black,thick] (0,0)--(-1,.5);
\draw[green!45!black,very thick] (-1,.5)--(1,.5);
\fill (0,0) circle(2pt) node[below left=7pt] {$O$};
\fill (1,.5) circle(2pt) node[above=14pt] {$A$};
\fill (-1,.5) circle(2pt) node[above=14pt] {$B$};
\end{scope}
\node at (3.55,2.6) {$AB:\ x=\frac12$，$k_{AB}$ 不存在};
\node at (10.65,2.6) {$AB:\ y=\frac12$，$k_{AB}=0$};
\node at (3.55,1.55) {$\displaystyle\frac{x_B-x_A}{y_B-y_A}=0$};
\node at (10.65,1.55) {$\displaystyle k_{OA}+k_{OB}=\frac12-\frac12=0$};
\node at (7.1,.55) {同一尺度下，横向模型的竖直弦关于 $y=x$ 翻转后成为纵向模型的水平弦。};
\end{tikzpicture}
\end{center}
\[
\boxed{\frac1{k_{OA}}+\frac1{k_{OB}}=\frac{x_B-x_A}{y_B-y_A}}.
\tag{H$^*$}
\]
当 $AB$ 不竖直时，右端就是 $1/k_{AB}$；当 $AB$ 竖直时，右端等于 $0$，但\textbf{绝不能称为 $1/k_{AB}$}（此时 $k_{AB}$ 不存在）。

对于 $x^2=2py$，不同点的横坐标必不同，因此 $AB$ 不可能竖直，公式 (V) 无须再排除竖直弦。

\textbf{特别提醒：}公式不依赖 $p$ 的正负。$p=0$ 不是本结论中的抛物线，不能代入。
\end{statementbox}

\clearpage
\begin{statementbox}
\textbf{直接推广：换一个顶点，或把顶点连线换成端点切线。}

\textbf{(1) 一般二次函数。}对 $y=ax^2+bx+c$，$a\ne0$，设顶点为 $V$，$A,B$ 为曲线上两个不同且均不为 $V$ 的点，则
\[
\boxed{k_{VA}+k_{VB}=k_{AB}}.
\]
平移不改变直线斜率，因此顶点不必在原点。这里 $a,b,c$ 为二次函数系数。

\medskip
\textbf{(2) 弦与两端切线。}此处 $A,B$ 只须是曲线上两个不同的点；$k_A,k_B$ 专指两点处的\textbf{切线斜率}，不是 $k_{OA},k_{OB}$。
对 $x^2=2py$，$p\ne0$，包括端点为顶点的情况，恒有
\[
\boxed{k_{AB}=\frac{k_A+k_B}{2}}.
\]
对 $y^2=2px$，$p\ne0$，须再要求 $y_Ay_B\ne0$ 且 $y_A+y_B\ne0$，即两端点不是顶点且弦不竖直，此时
\[
\boxed{\frac1{k_{AB}}=\frac12\left(\frac1{k_A}+\frac1{k_B}\right)}.
\]
横向顶点处的切线竖直，普通切线斜率不存在；横向竖直弦也不能写 $1/k_{AB}$。

\medskip
\textbf{快速变形：}横向顶点连线的斜率和非零时，
\[
\boxed{k_{AB}=\frac{k_{OA}k_{OB}}{k_{OA}+k_{OB}}}.
\]
看到倒数相加，先通分，再考虑“乘积除以和”；分母为零则改查竖直弦。

\textbf{关联提示：}后文另用焦点弦的\emph{长度倒数和}、直角三角形的\emph{长度平方倒数和}作对照；它们不是上述斜率公式的直接推广。
\end{statementbox}
